% Source-derived chapter 09: The main result
% Original source: chabauty-coleman-bound-surfaces
\chapter{The main result}
\label{chabauty-coleman-bound-surfaces-chapter-09}
\source{chabauty-coleman-bound-surfaces}

\begin{remark}[Source section]
\label{chabauty-coleman-bound-surfaces-chapter-09-source}
\source{chabauty-coleman-bound-surfaces}
\source{chabauty-coleman-bound-surfaces}
This chapter reproduces the corresponding section of the retrieved
LaTeX source, including its definitions, statements, examples, and
proofs. It has not yet been translated into Lean.
\notready
\end{remark}

% The source section title was: The main result
\label{SecMainResult}
\source{chabauty-coleman-bound-surfaces}

\subsection{Statement and first steps} We keep the notation from Section \ref{NotationLoc}. Let $\Acal/R$ be an abelian variety of relative dimension $n\ge 3$. Let $A=\Acal_K$ be the generic fibre and $A'=\Acal_\kk$ the special fibre. Let $\Xcal$ be an integral subscheme of $\Acal$ which is smooth, proper, of relative dimension $2$ over $R$ with geometrically irreducible fibres. Let $X=\Xcal_K$ and $X'=\Xcal_\kk$. Let $G\le A(K)$ be a finitely generated subgroup and let $\Gamma\subseteq A(K)$ be the $p$-adic closure of $G$ in $A(K)$. Our main result is:

%%
%%
\begin{theorem}[Main result]
\source{chabauty-coleman-bound-surfaces}
\notready\label{ThmMain} With the previous notation, suppose that $\rk (G)\le 1$, that 
\source{chabauty-coleman-bound-surfaces}
\begin{equation}\label{Eqnplarge}
\source{chabauty-coleman-bound-surfaces}
p>\max\left\{\efrak+1, \exp(\efrak/\exp(1))\right\} , 
\end{equation}
and that either of the following conditions holds:
\begin{itemize}
\item[(i)] $n=3$, $X$ is of general type, $X'$ contains no elliptic curves over $\kk^{alg}$, and 
$$
p> (128/9)\cdot c_1^2(X)^2.
$$ 
\item[(ii)] $A'$ is simple over $\kk^{alg}$ and there is an ample divisor $H$ on $A$ such that
$$
p> \max\left\{3 c_1^2(X)+2,\frac{n!\cdot (3\deg_H(X) + \cdeg_H(X))^n}{n^n\cdot \deg(H^n)}\right\}.
$$
\end{itemize}
Then $\Gamma\cap X(K)$ is finite and we have
\begin{equation}\label{EqnMainBound}
\source{chabauty-coleman-bound-surfaces}
\# \Gamma\cap X(K) \le \# X'(\kk) + \left(1-\frac{\efrak}{p-1}\right)^{-1}\cdot (q+4q^{1/2}+3)\cdot c_1^2(X).
\end{equation}
\end{theorem}
%%
%%



We keep the notation and assumptions of Theorem \ref{ThmMain} for the rest of this section. In addition we let $k=\kk^{alg}$. First we observe 

%%
%%
\begin{lemma}
\source{chabauty-coleman-bound-surfaces}
\notready\label{LemmaRk1} It suffices to prove Theorem \ref{ThmMain} under the assumption $\rk(G)=1$.
\source{chabauty-coleman-bound-surfaces}
\end{lemma}
\begin{proof} Suppose that $\rk(G)=0$ and let $\xi\in A(K)$ be a non-torsion point ($A(K)$ is uncountable). We may replace $G$ by the group generated by $G$ and $\xi$, which has rank $1$.
\end{proof}
%%
%%

Therefore, \emph{from now on we assume $\rk(G)=1$.}




The hypotheses of Theorem \ref{ThmMain} imply the following useful facts.

%%
%%
\begin{lemma}
\source{chabauty-coleman-bound-surfaces}
\notready\label{LemmaHyps} If either of (i) or (ii) in Theorem \ref{ThmMain} holds, then $X$ and $X'$ are surfaces of general type and $1\le c_1^2(X')=c_1^2(X)<(p-2)/3$. Furthermore:
\source{chabauty-coleman-bound-surfaces}

\begin{itemize}

\item If assumption (i) in Theorem \ref{ThmMain} holds, then $X'$  is an ample divisor on $A'$. 

\item If assumption (ii) in Theorem \ref{ThmMain} holds, then there is an ample divisor $H'$ on $A'$ such that $\deg_{H'}(X')=\deg_H(X)$, $\cdeg_{H'}(X')=\cdeg_H(X)$, and $\deg((H')^3)_{A'}=\deg(H^3)_A$.
\end{itemize}
\end{lemma}
\begin{proof} If $A'$ is simple over $k$ then $A$ is simple over $K^{alg}$ in which case $X\otimes K^{alg}$ does not contain translates of positive dimensional abelian subvarieties of $A\otimes K^{alg}$, and Lemma \ref{LemmaGenTypeTranslate} implies that $X$ is of general type (note that $K^{alg}\simeq \C$ as fields). Thus, if either of (i) or (ii) holds, $X$ is of general type and Theorem 9.1 in \cite{KatsuraUeno} implies that $X'$ is of general type. Lemma 9.3 in \cite{KatsuraUeno} gives $c_1^2(X')=c_1^2(X)$. We get $c_1^2(X)\ge 1$ by Lemma \ref{LemmaNef}. Either of (i) or (ii) implies $c_1^2(X)< (p-2)/3$.

If (i) holds then $X'$ is not an abelian surface because it is of general type, and since it does not contain elliptic curves over $k$, Lemma \ref{LemmaAmpleAbVar} implies that $X'$ is an ample divisor on $A'$. 

On the other hand, suppose that assumption (ii)  in Theorem \ref{ThmMain} holds. The  theory developed in the appendix to Expos\'e X in \cite{SGA6} (see also Section 20.3 in \cite{FultonInt}) gives a specialization map from the group of algebraic cycles  modulo numerical equivalence of $A$ to that of $A'$ which respects the intersection product. Let $H'$ be the specialization of $H$. Since $A'$ is geometrically simple, the divisor $H'$ is ample. Furthermore, since $\Xcal\to \Spec R$ is smooth, the specialization of a canonical divisor on $X$ is a canonical divisor on $X'$. The required equalities follow.
\end{proof}
%%
%%
\begin{remark}
\source{chabauty-coleman-bound-surfaces}
\notready\label{RmkMainToIntro} Theorem \ref{ThmMainIntro} follows from Theorem \ref{ThmMain}. Indeed, taking $K=\Q_p$ in Theorem \ref{ThmMain}, condition \eqref{Eqnplarge} becomes $p\ge 3$. Since $c_1^2(X)\ge 1$ (Lemma \ref{LemmaHyps}), either of (i) and (ii) implies $p\ge 7$, so the condition $p\ge 3$ can be dropped. Finally, \eqref{EqnMainBound} becomes \eqref{EqnMainIntro}.
\source{chabauty-coleman-bound-surfaces}
\end{remark}


We also have the following observation regarding curves in $X'_k$.

%%
%%
\begin{lemma}
\source{chabauty-coleman-bound-surfaces}
\notready\label{LemmaGenus2} Let $C$ be an irreducible curve in $X'_k$ defined over $k$. Then $\gfrak_g(C)\ge 2$.
\source{chabauty-coleman-bound-surfaces}
\end{lemma}
\begin{proof} Since $C\subseteq A'_k$ we have $\gfrak_g(C)\ge 1$. If it is $1$, then we get $\nu_{C/A'_k}:\widetilde{C}\to A'_k$ where $\widetilde{C}$ is an elliptic curve. Thus, $\nu_{C/A'_k}$ is, up to translation, a non-constant morphism of abelian varieties, and we get that its image $C$ is is an elliptic curve. But $X'_k$ contains no elliptic curves by assumption. 
\end{proof}
%%
%%

%%%%%%
%%%%%%
%%%%%%


\subsection{Choice of differentials and canonical divisor} \label{SecChoiceDiff} Recall that  we are assuming $\rk(G)=1$.
\source{chabauty-coleman-bound-surfaces}

%%
%%
\begin{lemma}
\source{chabauty-coleman-bound-surfaces}
\notready\label{LemmaChoice1PS} Let $\ttt$ be a choice of $R$-local parameters for $A$ at $e$. There is a $1$-parameter subgroup $\gamma=(\ttt,\uu)$ for $A$ such that $\Gamma\cap U_e\subseteq \im(\gamma)$. Moreover, $\gamma$ is uniquely determined, up to equivalence, by the condition that $\im(\gamma)\cap \Gamma$ strictly contains $e$.
\source{chabauty-coleman-bound-surfaces}
\end{lemma}
%%
%%
\begin{proof} By Lemma \ref{Lemma1PSgen} and the fact that $\rk(G)=1$,  there is a $1$-parameter subgroup $\gamma$ of the form $(\ttt,\uu)$ for the given $\ttt$, such that $\im(\gamma)\cap \Gamma$ strictly contains $\{e\}$. By Lemma \ref{LemmaSaturated} and the fact that $\Gamma$ is the $p$-adic closure of the rank $1$ group $G$, the non-triviality of $\im(\gamma)\cap \Gamma$ is equivalent to $\Gamma\cap U_e\subseteq \im(\gamma)$. By Lemma \ref{LemmaEquiv1PS}, the equivalence class of $\gamma$ is unique for this last property.
\end{proof}
%%
%%


In view of Lemma \ref{LemmaChoice1PS} and Lemma \ref{LemmaH1PS}, the group $G$ uniquely determines a $K$-linear hyperplane $\Hcal=\Hcal(\gamma)\subseteq H^0(A, \Omega^1_{A/K})$. 


%%
%%
\begin{lemma}
\source{chabauty-coleman-bound-surfaces}
\notready\label{LemmaDiff1} There are $\omega_1,\omega_2\in \Hcal$ such that
\source{chabauty-coleman-bound-surfaces}
\begin{itemize}
\item[(i)] $\omega_1,\omega_2$ are $K$-linearly independent.
\item[(ii)] $\omega_1,\omega_2\in H^0(\Acal,\Omega^1_{\Acal/R})$.
\item[(iii)] The differentials $\omega'_1,\omega'_2\in H^0(A',\Omega^1_{A'/\kk})$ obtained by reducing $\omega_1,\omega_2$ modulo $\pfrak$, are $\kk$-linearly independent.
\end{itemize}
\end{lemma}
\begin{proof} The ring $R$ is a DVR, hence $H^0(\Acal,\Omega^1_{\Acal/R})\simeq R^n$. Thus, the problem is reduced to showing the following: Given $n\ge 3$ and a $K$-linear map $f:K^n\to K$, there are $v_1,v_2\in \ker(f)\cap R^n$ which are linearly independent and such that their images in $\kk^n$ are $\kk$-linearly independent. For this, we may assume that $f$ is $R$-valued on $R^n$, and it follows that $\rk_R(\ker(f|_{R^n}))=n-1$. Thus, we can take $v_1,v_2\in \ker(f)\cap R^n$ which are $K$-linearly independent. After scaling, we can assume that $v_1,v_2\in R^n$ are primitive. If there are $b_1,b_2\in R^\times$ with $b_1 v_1+b_2 v_2\in \pfrak R^n\subseteq R^n$ then there are $c_1,c_2\in R$ such that $b_1 v_1+b_2 v_2=\varpi c_1 v_1+\varpi c_2v_2$ which is not possible since $v_1,v_2$ are $K$-linearly independent. Hence the images of  $v_1,v_2$ in $\kk^n$ are $\kk$-linearly independent.
\end{proof}
%%
%%

From now on we fix a choice of $\omega_1,\omega_2$ as in Lemma \ref{LemmaDiff1} and let $\omega'_1,\omega'_2\in H^0(A',\Omega^1_{A'/\kk})$ be their reductions modulo $\pfrak$. Let $\iota:X'\to A'$ be the inclusion map. For $j=1,2$ we let $u_j=\iota^\bullet(\omega'_j)\in H^0(X',\Omega^1_{X'/\kk})$ be the restriction of $\omega'_j$ to $X'$.

%%
%%
\begin{lemma}
\source{chabauty-coleman-bound-surfaces}
\notready\label{LemmaDiff2} The $2$-form   $u_1\wedge u_2\in H^0(X',\Omega^2_{X/\kk})$ is not the zero section.
\source{chabauty-coleman-bound-surfaces}
\end{lemma}
\begin{proof} Since $A'$ is an abelian variety,  (iii) in Lemma \ref{LemmaDiff1} shows that $\omega'_1\wedge \omega'_2\in H^0(A',\Omega^2_{A'/\kk})$ is not the zero section. By the assumptions (i) and (ii) in Theorem \ref{ThmMain} and by Lemma \ref{LemmaHyps}, we can apply Lemma \ref{Lemma2forms} to $\omega'_1,\omega'_2$ on $A'$ and the surface $X'\subseteq A'$ after base change to $k$. 
\end{proof}
%%
%%

 By Lemma \ref{LemmaDiff2} we can define $D=\divi_{X'}(u_1\wedge u_2)$ on $X'$. Lemmas \ref{LemmaHyps}, \ref{LemmaGenus2},  and \ref{LemmaGTA} give:
 
 \begin{lemma}
\source{chabauty-coleman-bound-surfaces}
\notready\label{LemmaAAA}  $D$ is effective, ample,  and it is a canonical divisor of $X'$ defined over $\kk$. In particular, $D$ is numerically effective and $D> 0$.
\source{chabauty-coleman-bound-surfaces}
\end{lemma}


Let $Z=\supp(D)$ and let $C_1,...,C_\ell$ be the irreducible components of $Z$ over $k$. Thus, $Z$ is a (possibly singular and reducible) curve defined over $\kk$ and we have $D=\sum_{j=1}^\ell a_j C_j$ as divisors on $X'_k=X'\otimes k$ for certain integers $a_j\ge 1$. Let us write $\nu_j=\nu_{C_j/X'_k}:\widetilde{C}_j\to X'_k$ for $j=1,...,\ell$.

%%
%%
\begin{lemma}
\source{chabauty-coleman-bound-surfaces}
\notready\label{LemmaChoicew} There is $b\in \kk$ such that the differential $w=u_1+bu_2\in H^0(X',\Omega^1_{X'/\kk})$ satisfies that no map $\nu_j:\widetilde{C}_j\to X'_k$ is $w$-integral.
\source{chabauty-coleman-bound-surfaces}
\end{lemma}
\begin{proof} By Lemmas \ref{LemmaHyps} and \ref{LemmaAAA} (especially, ampleness of $D$) we have 
$$
\ell \le \sum_{j=1}^\ell a_j \le \sum_{j=1}^\ell a_j\cdot (C_j.D)_{X'_k}=c_1^2(X')<p
$$
Given $b\in \kk$ let $w_b=u_1+bu_2\in H^0(X',\Omega^1_{X'/\kk})$. If for every $b\in \F_p\subseteq \kk$ we have that some $\nu_{j}$ is $w_b$-integral, the fact that $\ell<p$ implies that for some $j_0$ there are $b\ne b'$ in $\kk$ such that $w_b,w_{b'}\in \ker(\nu_{j_0}^\bullet)$. Hence, $\omega_1+b\omega_2$ and $\omega_1+b'\omega_2$ are in $\ker (\nu_{C_j/A'_k}^\bullet)$. Since $b\ne b'$ we deduce $\dim_k \ker (\nu_{C_j/A'_k}^\bullet)\ge 2$. By the assumptions (i) and (ii) in Theorem \ref{ThmMain} and Lemma \ref{LemmaHyps}, this contradicts Lemma \ref{LemmaSemiInj}.
\end{proof}
%%
%%

From now on, we fix a choice of $b\in \kk$ and the corresponding $w=u_1+bu_2\in H^0(X',\Omega^1_{X'/\kk})$ as in Lemma \ref{LemmaChoicew}.



%%%%%%
%%%%%%
%%%%%%


\subsection{Bounds from overdetermined $\omega$-integrality}\label{SecBoundsOver} %
\source{chabauty-coleman-bound-surfaces}

Given $x\in X'(k)$ we let $m_{X',\omega'_1,\omega'_2}(x)$ be the supremum (possibly infinite) of all integers $m\ge 0$ with the property that there is a closed immersion $\phi : V_m^k\to X'_k$ supported at $x$ which is $\omega$-integral for both $\omega=u_1,u_2$. Here we recall from Section \ref{SecOver} that $V^k_m=\Spec k[z]/(z^{m+1})$. We observe that, fixing the ambient abelian variety $A'$ over $\kk$,  the quantity $m_{X',\omega'_1,\omega'_2}(x)$ only depends on our choices of $X'$, $\omega'_1$, $\omega'_2$, and $x$. As $X'$ is given  and $\omega'_1,\omega'_2$ are fixed, we may simply write $m(x)=m_{X',\omega'_1,\omega'_2}(x)$ (except in the proof of Lemma \ref{LemmaXY} below, where a reduction argument will require to consider other choices of $X'$).

%%
%%
\begin{lemma}
\source{chabauty-coleman-bound-surfaces}
\notready\label{Lemmamx0} If $x\notin Z$ then $m(x)=0$.
\source{chabauty-coleman-bound-surfaces}
\end{lemma}
\begin{proof} By Theorem \ref{ThmOver}. See Remark \ref{RmkOver1} for details in the case $x\notin Z=\supp(D)$.
\end{proof}
%%
%%

%%
%%
\begin{lemma}
\source{chabauty-coleman-bound-surfaces}
\notready\label{Lemmamx1} If $x\in Z$ then 
\source{chabauty-coleman-bound-surfaces}
$$
m(x)\le \sum_{j=1}^\ell \sum_{y\in \nu_j^{-1}(x)} a_j\cdot (\ord_y(\nu_j^\bullet(w))+1).
$$
\end{lemma}
\begin{proof} By Theorem \ref{ThmOver}, after choosing $\omega_0=w=u_1+bu_2$.
\end{proof}
%%
%%




%%
%%
\begin{lemma}
\source{chabauty-coleman-bound-surfaces}
\notready\label{Lemmamxp} For every $x\in X'(k)$ we have $m(x)\le p-3$. In particular, $m(x)$ is finite.
\source{chabauty-coleman-bound-surfaces}
\end{lemma}
\begin{proof} By Lemma \ref{Lemmamx0} we may assume $x\in Z$. Lemma \ref{Lemmamx1} gives
$$
m(x)\le \sum_{j=1}^\ell  a_j\cdot \deg_{\widetilde{C}_j}(\nu_j^\bullet(w)) + \sum_{j=1}^\ell a_j\cdot \# \nu_j^{-1}(x).
$$
Since $w$ is chosen as in Lemma \ref{LemmaChoicew} we get $\deg_{\widetilde{C}_j}(\nu_j^\bullet(w))=2\gfrak_g(C_j)-2\le 2\gfrak_a(C_j)-2$. On the other hand, Lemma \ref{LemmaGenus2} gives $\gfrak_g(C_j)\ge 2$ for each $j$, and in view of Lemmas \ref{LemmaHironakaGenus} and \ref{LemmaBounddelta} we get
$$
\#\nu^{-1}_j(x)\le \delta(C_j,x)+1\le \gfrak_a(C_j)-1.
$$
Using Lemmas \ref{LemmaCanonicalGenusBd}, \ref{LemmaHyps}, and \ref{LemmaAAA} this gives $m(x)\le 3\cdot  \sum_{j=1}^\ell a_j\cdot (\gfrak_a(C_j)-1)\le 3c_1^2(X')<p-2$.
\end{proof}
%%
%%



%%%%%%
%%%%%%
%%%%%%


\subsection{Bounds on residue disks}\label{SecResDisk} %
\source{chabauty-coleman-bound-surfaces}



Every $\xi\in A(K)$ extends to a section $\sigma_\xi:\Spec R\to \Acal$ because $\Acal\to \Spec R$ is proper. This determines a reduction map $\red: A(K)\to A'(\kk)$ which is a group morphism. For $x\in A'(\kk)$ we let $U_x=\red^{-1}(x)\subseteq A(K)$ be the corresponding residue disk. In particular, $U_e=\ker(\red)$ as in Section \ref{SecLocLin}. The goal of this paragraph is to show the following result (under the assumptions of Theorem \ref{ThmMain}), which can be of independent interest:

%%
%%
\begin{proposition}
\source{chabauty-coleman-bound-surfaces}
\notready\label{PropKeyU} Let $\xi\in \Gamma\cap X(K)$ and  $x=\red(\xi)\in X'(\kk)$. Then $\Gamma\cap X(K)\cap U_x$ is finite and
\source{chabauty-coleman-bound-surfaces}
$$
\# \Gamma\cap X(K)\cap U_x\le 1+ m(x)\cdot  \left(1-\frac{\efrak}{p-1}\right)^{-1}.
$$
In particular, if $m(x)=0$, then $U_x$ contains at most one point of $\Gamma\cap X(K)$.
\end{proposition}
%%
%%

Recall that we are assuming $\rk(G)=1$. In addition, we can make the following reduction in the proof of Proposition \ref{PropKeyU}.

%%
%%
\begin{lemma}
\source{chabauty-coleman-bound-surfaces}
\notready\label{LemmaXY} It suffices to prove Proposition \ref{PropKeyU} under the assumption $\xi=e$.
\source{chabauty-coleman-bound-surfaces}
\end{lemma} 
\begin{proof} Let $\xi_0\in \Gamma\cap X(K)$ and consider the translation $\Ycal=\Xcal- \sigma_{\xi_0}(\Spec R)$. Note that the generic and special fibres of $\Ycal$ are $Y=X-\xi_0$ and $Y'=X'-x_0$ where $x_0=\red(\xi_0)$. The assumptions of Theorem \ref{ThmMain} also hold for $\Ycal$ instead of $\Xcal$. Since $\Gamma-\xi_0=\Gamma$, we find 
$$
\# \Gamma\cap X(K)\cap U_{x_0} = \# \Gamma \cap Y(K)\cap U_e.
$$ 
On the other hand, as the differentials $\omega'_1,\omega'_2\in H^0(A',\Omega^1_{A'/\kk})$ are translation invariant, we have 
$$
m_{X',\omega'_1,\omega'_2}(x_0)=m_{Y',\tau^*\omega'_1,\tau^*\omega'_2}(e)=m_{Y',\omega'_1,\omega'_2}(e)
$$ 
where $\tau:A'\to A'$ is the translation map $P\mapsto P+x_0$. The differential $\omega'_j$ is determined by $\omega_j\in H^0(A, \Omega^1_{A/K})$ (for $j=1,2$), which in turn is chosen according to Lemmas \ref{LemmaChoice1PS} and  \ref{LemmaDiff1}. Thus, the choice of $\omega'_1,\omega'_2$  only depends on $\Acal$ and $G$, not on the particular choice of $\Xcal$. Therefore, if Proposition \ref{PropKeyU} holds in the case $\xi=e$ for the given $\Acal$ and $G$, and for every choice of $\Xcal$ as in Theorem \ref{ThmMain}, then we can in particular apply it to $\Ycal$ obtaining 
$$
\# \Gamma\cap X(K)\cap U_{x_0} = \# \Gamma \cap Y(K)\cap U_e\le m_{Y',\omega'_1,\omega'_2}(e)=m_{X',\omega'_1,\omega'_2}(x_0).
$$
\end{proof}
%%
%%

For the rest of the current Section \ref{SecResDisk}, let us assume that $e\in \Gamma\cap X(K)$ in order to show Proposition \ref{PropKeyU} for $\xi=e$. This is enough, by Lemma \ref{LemmaXY}.

Let $t_1,...,t_n\in \mfrak_{\Acal, e}$ be a system of $R$-local parameters for $\Acal$ along $\sigma_e$ with the property that $t_1, ..., t_{n-2}$ are a system of local equations for $\Xcal$ along $\sigma_e$. This is possible by Lemma \ref{LemmaChoiceLocParam}. Let $t'_j$ be the restriction of $t_j$ to $A'$. Then $t'_1,...t'_{n-2}$ is a system of local equations for $X'\subseteq A'$ at $e$. In particular, $t_1,...,t_{n-2}$ restricted to $A$ vanish on $X(K)\cap U_e$.


Write $\ttt=(t_1,...,t_n)$ and let $\uu\in K^n$ with $|\uu |=1$ be such that $\gamma=(\ttt,\uu)$ be a $1$-parameter subgroup of $A$ with $\Gamma\cup U_e\subseteq \im(\gamma)$. By Lemma \ref{LemmaChoice1PS}, the vector $\uu$  exists, $\gamma$ is unique up to equivalence, and $\gamma$ corresponds to the hyperplane $\Hcal\subseteq H^0(A,\Omega^1_{A/K})$ from Section \ref{SecChoiceDiff}. We get
$$
\# \Gamma\cap X(K)\cap U_e \le \# X(K)\cap \im(\gamma).
$$

By Lemma \ref{LemmaKeyLogExp} (which is applicable by \eqref{Eqnplarge}), the definition of the map $\tilde{\gamma}:\pfrak\mapsto U_e$ (cf. \eqref{Eqntildegamma}), and the fact that $t_1,...,t_{n-2}$ vanish on $X(K)\cap U_e$, we deduce that for each $j=1,...,n-2$
$$
 \# X(K)\cap \im(\gamma)\le \#\{z_0\in \pfrak : \Exp_j^\ttt(z_0\cdot \uu)=0\}.
$$
Therefore, for each $j=1,..., n-2$ we obtain
%%
\begin{equation}\label{EqnBdUn0}
\source{chabauty-coleman-bound-surfaces}
\# \Gamma\cap X(K)\cap U_e \le \nn_0(\Exp_j^\ttt(z\cdot \uu), 1/q).
\end{equation}
%%
Let $P^\ttt_{j,h}(\xx)\in K[x_1,...,x_n]$ be the homogeneous part of degree $h$ in $\Exp_j^\ttt(\xx)$. Note that $P^\ttt_{j,0}(\xx)=0$.

%%
%%
\begin{lemma}
\source{chabauty-coleman-bound-surfaces}
\notready\label{LemmaBdNm} There are positive integers integers  $j_0\le n-2$ and $N\le m(e)+1$ with $|P^{\ttt}_{j_0,N}(\uu)|\ge 1$.
\source{chabauty-coleman-bound-surfaces}
\end{lemma}
\begin{proof} Let $0\le m_0\le p-2$ be an integer satisfying $|P^{\ttt}_{j,h}(\uu)|< 1$ for each $0\le h\le m_0$ and for each $1\le j\le n-2$. We claim that $m_0\le m(e)$.

Since $\omega_1,\omega_2$ are chosen as in Lemma \ref{LemmaDiff1}, we see that Lemmas \ref{LemmaClosedImm} and \ref{Lemmawintgamma} give a closed immersion $\tilde{\gamma}'_{m_0}: V'_m\to A'$ which is $\omega'_i$-integral for $i=1,2$. Let us prove that $\tilde{\gamma}'_{m_0}$ factors through a closed immersion $\phi_{m_0}:V'_{m_0}\to X'$ supported at $x=e$. Indeed, let $G_{j,m_0}(z)\in K[z]$ be the truncation of $\Exp^\ttt_{j}(z\cdot \uu)$ up to degree $m_0$, that is,
$$
G_{j,m}(z)=\sum_{h=0}^m  P^{\ttt}_{j,h}(\uu)\cdot z^h\in K[z].
$$
Our assumption $|P^{\ttt}_{j,h}(\uu)|< 1$ for all $0\le h\le m_0$ and $1\le j\le n-2$ implies that $G_{j,m_0}(z)\in R[z]$ for all $1\le j\le n-2$, and in fact, in this case $G_{j,m_0}(z)\equiv 0 \bmod \pfrak$. In view of Lemma \ref{LemmaTruncate1ps} which gives the construction of $\tilde{\gamma}'_{m_0}$, reducing  \eqref{Eqninf1psP} modulo $\pfrak$ we get  that for each $1\le j\le n-2$
$$
(\tilde{\gamma}'_{m_0})^\#(t'_j)= G_{j,m_0}(z) \bmod (z^{m_0+1},\varpi)=0.
$$
As $t'_1,...,t'_{n-2}$ are local equations for $X'$ at $e$, this shows that $\tilde{\gamma}'_{m_0}$ factors through the inclusion $\iota:X'\to A'$. We obtain the claimed closed immersion $\phi_{m_0}:V'_{m_0}\to X'$ satisfying $\iota\circ \phi_{m_0}=\tilde{\gamma}'_{m_0}$.

Since $\tilde{\gamma}'_{m_0}$ is $\omega'_i$-integral for $i=1,2$, we obtain that $\phi_{m_0}$ is $u_i$-integral for $i=1,2$. The same holds for the base change of $\phi_{m_0}$ to $k$, which is a closed immersion $V^k_{m_0}\to X'_k$ supported at $x=e$. Thus, by definition of $m_{X',\omega'_1,\omega'_2}(x)$ (cf. Section \ref{SecBoundsOver}) we finally get $m_0\le m(e)$, as claimed.

Suppose now that $|P^{\ttt}_{j,h}(\uu)|< 1$ for each $0\le h\le p-2$ and $1\le j\le n-2$. Taking $m_0=p-2$ in our initial claim we would get $p-2\le m(e)$, which is not possible by Lemma \ref{Lemmamxp}. 

Therefore, there is some $0\le h_0\le p-2$ and some $1\le j_0\le n-2$ such that $|P^{\ttt}_{j_0,h_0}(\uu)|\ge 1$. Let $N$ be the least value of such an $h_0$ and choose any $1\le j_0\le n-2$ satisfying $|P^{\ttt}_{j_0,N}(\uu)|\ge 1$. Thus, $N$ exists and $N\le p-2$. Also, $N\ge 1$ since $P^{\ttt}_{j,0}(\uu)=0$ for each $j$. Taking $m_0=N-1\le p-3$ we see that $|P^{\ttt}_{j,h}(\uu)|< 1$ for each $h\le m_0$ and  $j\le n-2$, and out initial claim implies $N-1=m_0\le m(e)$.
\end{proof}
%%
%%


\begin{proof}[Proof of Proposition \ref{PropKeyU}] As proved in Lemma \ref{LemmaXY}, we may assume that $e\in \Gamma\cap X(K)$ and that $\xi=e$. Let $j_0$ and $N$ be as in Lemma \ref{LemmaBdNm}, let $M=p^{[K:\Q_p]/(p-1)}$, and let $r=1/q$. Let us write $\lambda=\efrak/(p-1)$ and observe that $\lambda<1$ by \eqref{Eqnplarge}. We remark that $M>1$, and $r<M^{-1}$ because
%%
\begin{equation}\label{EqnFracMr}
\source{chabauty-coleman-bound-surfaces}
\frac{\log M}{\log r^{-1}} =\frac{\log M}{\log q} = \frac{[K:\Q_p]}{(p-1)\ffrak} = \lambda<1.
\end{equation}
%%
If $c_{j,\alpha}\in K$ is the coefficient of $\xx^\alpha$ in $\Exp^\ttt_{j}(\xx)$, Lemma \ref{LemmaCvExp} gives $|c_{j,\alpha}| \le M^{\| \alpha\|-1}$. Hence, we can apply Lemma \ref{LemmaMVzeros} with these choices to the power series $\Exp^\ttt_{j_0}(\xx)$ obtaining
$$
\nn_0(\Exp_{j_0}^\ttt(z\cdot \uu), 1/q)\le \left(N-\lambda\right)\left(1-\lambda\right)^{-1}
$$
where we used \eqref{EqnFracMr}. By \eqref{EqnBdUn0} and Lemma \ref{LemmaBdNm} we conclude
$$
\# \Gamma\cap X(K)\cap U_e \le  \left(m(e)+1-\lambda\right)\left(1-\lambda\right)^{-1}=1+ m(e)\cdot \left(1-\lambda\right)^{-1}.
$$
\end{proof}

%%%%%%
%%%%%%
%%%%%%


\subsection{Adding over residue disks} 

\begin{proof}[Proof of Theorem \ref{ThmMain}] By Proposition \ref{PropKeyU} and  Lemma \ref{Lemmamx0} we obtain
%%
\begin{equation}\label{EqnFinalU}
\source{chabauty-coleman-bound-surfaces}
\begin{aligned}
\# \Gamma\cap X(K)& =\sum_{x\in X'(\kk)} \# \Gamma \cap X(K)\cap U_x= \# X'(\kk) + \left(1-\frac{\efrak}{p-1}\right)^{-1}\sum_{x\in X'(\kk)} m(x)\\
&= \# X'(\kk) + \left(1-\frac{\efrak}{p-1}\right)^{-1}\sum_{x\in Z(\kk)} m(x).
\end{aligned}
\end{equation}
%%
Recall from Lemma \ref{LemmaChoicew} that for each $j=1,...,\ell$ we have that $\nu_j^\bullet(w)$ is not the zero differential. By Lemma \ref{Lemmamx1} we have have 
%%
\begin{equation}\label{EqnFinalm}
\source{chabauty-coleman-bound-surfaces}
\begin{aligned}
\sum_{x\in Z(\kk)} m(x)&\le \sum_{x\in Z(\kk)} \sum_{j=1}^\ell \sum_{y\in \nu_j^{-1}(x)} a_j\cdot (\ord_y(\nu_j^\bullet(w))+1)\\
&\le \sum_{j=1}^\ell a_j\cdot \deg_{\widetilde{C}_j}(\nu_j^\bullet (w)) + \sum_{x\in Z(\kk)} \sum_{j=1}^\ell a_j\cdot \#\nu_j^{-1}(x).
\end{aligned}
\end{equation}
%%
Lemma \ref{LemmaCanonicalGenusBd} (which is applicable by Lemma \ref{LemmaAAA}) and Lemma  \ref{LemmaHyps} give
%%
\begin{equation}\label{EqnFinalS1}
\source{chabauty-coleman-bound-surfaces}
\sum_{j=1}^\ell a_j\cdot \deg_{\widetilde{C}_j}(\nu_j^\bullet (w)) = \sum_{j=1}^\ell a_j\cdot (2 \gfrak_g({C}_j) -2)\le 2c_1^2(X')=2c_1^2(X).
\end{equation}
%%
Let $D_1,..., D_s$ be the irreducible components of $Z$ over $\kk$ and let $J_1,...,J_s$ be the partition of $\{1,2,...,\ell\}$ determined by the condition $D_i\otimes k=\cup_{j\in J_i} C_j$. Since each $D_i$ is defined over $\kk$, there are integers $b_i\ge 0$ for $1\le i\le s$ such that $a_j=b_i$ for each $j\in J_i$. By Lemma \ref{LemmaWeilBound}   we obtain
$$
\begin{aligned}
\sum_{x\in Z(\kk)} \sum_{j=1}^\ell a_j\cdot \#\nu_j^{-1}(x) & =\sum_{i=1}^s b_i \cdot \sum_{x\in D_i(\kk)} \sum_{j\in J_i} \#\nu_j^{-1}(x) \\
& \le \sum_{i=1}^s b_i \cdot \left(  (q+1)\cdot \#J_i + 2q^{1/2}\cdot \sum_{j\in J_i} \gfrak_g(C_j) \right)\\
&= (q+1)\sum_{j=1}^\ell a_j + 2q^{1/2}\cdot \sum_{j=1}^\ell a_j\cdot \gfrak_g(C_j)\\
&= (q+2q^{1/2}+1)\sum_{j=1}^\ell a_j + 2q^{1/2}\cdot \sum_{j=1}^\ell a_j\cdot (\gfrak_g(C_j)-1).
\end{aligned}
$$
From Lemmas \ref{LemmaGenus2} and \ref{LemmaCanonicalGenusBd} (applicable by Lemma \ref{LemmaAAA}) we deduce
%%
\begin{equation}\label{EqnFinalS2}
\source{chabauty-coleman-bound-surfaces}
\sum_{x\in Z(\kk)} \sum_{j=1}^\ell a_j\cdot \#\nu_j^{-1}(x) \le (q+4q^{1/2}+1)\sum_{j=1}^\ell a_j\cdot (\gfrak_g(C_j)-1)\le (q+4q^{1/2}+1)c_1^2(X).
\end{equation}
%%
Using \eqref{EqnFinalS1} and \eqref{EqnFinalS2} in \eqref{EqnFinalm} we get
$$
\sum_{x\in Z(\kk)} m(x)\le (q+4q^{1/2}+3)c_1^2(X)
$$
which together with \eqref{EqnFinalU} gives the result.
\end{proof}


%%%%%%%%%%%%%%%%%%%%%%%%%%%%%%%%%%%%%%
%%%%%%%%%%%%%%%%%%%%%%%%%%%%%%%%%%%%%%
%%%%%%%%%%%%%%%%%%%%%%%%%%%%%%%%%%%%%%
%%%%%%%%%%%%%%%%%%%%%%%%%%%%%%%%%%%%%%
%%%%%%%%%%%%%%%%%%%%%%%%%%%%%%%%%%%%%%
%%%%%%%%%%%%%%%%%%%%%%%%%%%%%%%%%%%%%%
